\section{Key Terminology}\label{sec:introduction:terminology}
The following terms are used with specific meaning throughout this thesis.
Standard quantum computing vocabulary is defined in Section~\ref{sec:introduction:quantum_computing_foundations}; standard programming language theory terms are defined at point of first use in the relevant chapters.

\subsection{Introduced terms}\label{subsec:introduction:terminology:introduced_terms}

\begin{itemize}
    \item \textbf{Quantum-implicit.} A class of quantum programming language whose primary surface model is expressed in ordinary mathematical and computational abstractions.
    A quantum-implicit language allows programmers to describe and implement quantum algorithms without explicitly annotating or reasoning about qubits, gates, circuits or other quantum-mechanical constructs.
    The compiler is responsible for inferring and synthesising the necessary quantum realisation.

    \item \textbf{Quantum bias.} The degree to which a language's surface model requires the programmer to understand quantum-mechanical vocabulary to write correct programs.
    A language with high quantum bias exposes circuit-level concepts as mandatory programmer-facing constructs.
    A language with zero quantum bias is quantum-implicit.
    The concept is formally operationalised in Chapter~\ref{ch:pub_qpl_review}.
\end{itemize}

\subsection{Borrowed terms}\label{subsec:introduction:terminology:borrowed_terms}

\begin{itemize}
    \item \textbf{Surface model.} The vocabulary, abstractions and mental model a programming language presents to the programmer, as distinct from the language's compilation target, implementation or runtime.
    In this thesis, the surface model is the primary object of analysis: what the surface model requires of the programmer, not what the language generates at the hardware level.

    \item \textbf{Bleed-through.} The leakage of execution-layer or hardware-layer concerns into the programming-layer surface model~\cite{DiMatteo2024AbstractionHierarchy}.
    A language exhibits bleed-through when the programmer must reason about circuit depth, qubit allocation, gate set constraints or decoherence timescales to write correct programs.
    Bleed-through is the mechanism that produces quantum bias.

    \item \textbf{Quantum substrate.} The underlying physical quantum-mechanical model: qubits, unitary gate operations, measurement and the decoherence constraints of real hardware.
    A quantum-implicit language is agnostic of the quantum substrate; its surface model does not expose substrate concepts, even though the compiled output must address them.

    \item \textbf{Gate-based quantum hardware.} Quantum hardware that executes programs as sequences of discrete unitary gate operations on qubit registers.
    This is the predominant hardware model addressed by quantum programming languages and the compilation target for all four incorporated publications.
    It is distinct from quantum annealing hardware and continuous-variable quantum systems.
\end{itemize}

\subsection{Abbreviations}\label{subsec:introduction:terminology:abbreviations}

\begin{itemize}
    \item \textbf{NISQ (Noisy Intermediate-Scale Quantum).} The current generation of quantum hardware, characterised by qubit counts in the range of tens to hundreds,without full fault-tolerant error correction.
    All empirical results in this thesis are evaluated against NISQ-class hardware.

    \item \textbf{DSR (Design Science Research).} The research methodology adopted in this thesis.
    Full treatment is in Chapter~\ref{ch:methodology}.

    \item \textbf{QPL (Quantum Programming Language).} Used as shorthand for any quantum programming language, framework or embedded domain-specific language that provides a programmer-facing interface for expressing quantum computations.
\end{itemize}